\clearpage
\section{Language modeling as a common task}\label{sec:task}
\subsection*{From text applications to a probability model}
A translation, a summary, and a continuation all contain sequences of text.
A language model assigns probabilities to token sequences. This makes it
possible to score candidate sequences and generate new ones. A high model
probability expresses what the model expects; factual accuracy and success
on a task still need their own checks.

Let $x_1,\ldots,x_T$ be tokens from a finite output vocabulary
$\mathcal V$ of size $K$. We first fix the tokenizer and the treatment of
document boundaries. The chain rule gives the exact identity
\[
 p_\theta(x_{1:T})=\prod_{t=1}^{T}p_\theta(x_t\mid x_{<t}).
\]
This identity makes no short-context assumption. An $n$-gram, a feedforward
network, and an RNN differ in how they approximate each conditional.
Every conditional must be nonnegative and sum to one over $\mathcal V$.

\subsection*{The learning objective}
Write $P$ for the data distribution and $Q_\theta$ for the model. On a
fixed sequence sample space, with $Q_\theta$ positive wherever $P$ is positive,
\[
 D_{\mathrm{KL}}(P\Vert Q_\theta)
 =\mathbb E_P[\log P(X)]-\mathbb E_P[\log Q_\theta(X)].
\]
The first term is independent of $\theta$, so minimizing this divergence
amounts to maximizing expected log likelihood. We cannot sum over the
unknown data distribution directly. Instead, training uses observed documents
and minimizes their empirical negative log likelihood:
\[
 \mathcal L(\theta)=-\frac{1}{N_{\mathrm{pred}}}
 \sum_{d}\sum_{t\in\mathcal S_d}\log p_\theta(x_t^{(d)}\mid x_{<t}^{(d)}).
\]
Here $\mathcal S_d$ specifies the scored positions, including an end token
if the evaluation convention scores it. $N_{\mathrm{pred}}$ is their total
count. During training this denominator is constant with respect to the
parameters. Averaging each document separately and then averaging documents
would give shorter documents more weight per token.

\subsection*{Worked example: probability is a product, loss is a sum}
Suppose the model assigns probabilities $1/2$, $1/4$, and $1/2$ to the
three scored events in a short document. The document probability is $1/16$,
its negative log probability is $\log 16$, and its mean token loss is
$\log(16)/3\approx0.9242$ nats. Taking logs converts a long product of
small numbers into a sum and makes the contribution of each prediction visible.

\textbf{Boundaries matter.} A beginning marker supplies context but need
not be an output symbol. An end marker can be a predicted event. Including
the probability of termination is part of defining a distribution over
variable-length documents; a generation length cap is a separate practical rule.

\takeaway{Fix the representation and scoring events before comparing models.
Then ask how each model turns the allowed context into a normalized distribution.}
\selfcheck{If one target receives probability zero, what happens to the document
probability and its loss? The answers are zero and $+\infty$. Explain why
adding an end token changes the scored-event count.}
\textbf{Reading connection.} \readingref{R01} connects uncertainty with
coding; \readingref{R04} replaces discrete context tables with shared learned features.
